\documentclass[TFG_spanish.tex]{subfiles}

\begin{document}

\clearpage
\thispagestyle{empty}

\begin{center}
    \large \textbf{Declaración de originalidad}
\end{center}

{\fontsize{11}{15}\selectfont

    Declaro bajo mi responsabilidad que el Proyecto presentado con el título \textbf{\projectTitle} en la Escuela Técnica Superior de Ingeniería ICAI de la Universidad Pontificia Comillas en el curso académico \BeginYear/\EndYear\ es de mi autoría y no ha sido presentado anteriormente para otros fines. El Proyecto no ha sido plagiado de ningún otro, ni total ni parcialmente, y la información que ha sido tomada de otros documentos está debidamente referenciada.

    \vspace{5mm}

    \begin{center}
        {\large \textbf{Uso de Inteligencia Artificial}}\footnote{Esta declaración se refiere al uso de Inteligencia Artificial generativa para la elaboración de los documentos del Proyecto (Anexo B y Memoria). No se aplica a Proyectos en los que, por su naturaleza, deba utilizarse inteligencia artificial como parte de los mismos (aplicación de técnicas de aprendizaje automático, redes neuronales, análisis de datos...).}
    \end{center}

     Declaro bajo mi responsabilidad (indicar la opción correcta):

    \begin{itemize}
        \item[$\square$] % Sustituir \square por \boxtimes para seleccionar esta opción
        No he utilizado Inteligencia Artificial en la elaboración de este documento.

        \item[$\boxtimes$] % Opción seleccionada
        He utilizado Inteligencia Artificial en la elaboración de este documento y/o del Anexo B bajo las condiciones permitidas por la Universidad Pontificia Comillas, es decir, aplicando el Nivel 2 de la \href{https://aiassessmentscale.com/}{Escala de Evaluación de Perkins et al.\ (2024)}: \emph{``La IA puede utilizarse para actividades previas a la tarea, como lluvia de ideas, descripción e investigación inicial. Este nivel se centra en el uso de la IA para planificar, sintetizar y generar ideas, pero las evaluaciones deben enfatizar la capacidad de desarrollar y perfeccionar estas ideas de forma independiente''}. En concreto, la Inteligencia Artificial se ha utilizado para:
    \end{itemize}

    \begin{center}
    \renewcommand{\arraystretch}{3}
    \begin{tabular}{|p{\textwidth}|}
    \hline
    Se ha utilizado IA generativa (Nivel 2) como apoyo para comprender conceptos metodológicos y estadísticos, mejorar la redacción y la coherencia del texto, y verificar referencias. El diseño, la implementación, los experimentos y las decisiones del trabajo son obra propia del autor.
    \vspace{4cm}\\
    \hline
    \end{tabular}
    \end{center}

    \vspace{1.5em}

    \begin{center}
    \renewcommand{\arraystretch}{2}
    \begin{tabular}{|p{0.9\textwidth}|}
    \hline
    \begin{center}
    \vspace{5mm}
        (firmar aquí)
    \end{center}
    \\
    \hline
    Firma: \Author \\
    \hline
    Fecha: 13/06/\EndYear\\
    \hline
    \end{tabular}
    \end{center}

    \begin{center}
        \textbf{Autorización para la entrega del Proyecto}
    \end{center}

    \begin{center}
    \renewcommand{\arraystretch}{2}
    \begin{tabular}{|p{0.45\textwidth} |p{0.45\textwidth} |
    }
        \hline
        \begin{center} \vspace{-8mm} Director del TFG \end{center} &
        \begin{center} \vspace{-8mm} Codirector del TFG, en su caso \end{center}\\
        \hline
        \begin{center}
            \vspace{5mm}
            (firmar aquí)
        \end{center} &
        \begin{center}
            \vspace{5mm}
            (firmar aquí)
        \end{center}
        \\
        \hline
        Firma: \Director & Firma: \CoDirector \\
        \hline
        Fecha: \hspace{0.9cm}/\hspace{0.9cm}/\EndYear
        &
        Fecha: \hspace{0.9cm}/\hspace{0.9cm}/\EndYear\\
        \hline
    \end{tabular}
    \end{center}
}



\clearpage

\end{document}
